
\assignmenttitle
    {LAB ASSIGNMENT 5}
    {INHERITANCE}

\aim{
    To write C++ programs utilizing inheritance, polymorphism, constructors,
    access specifiers, and virtual functions.
}

\programs

\program{
    1. Create a base class Employee with attributes name and salary, and a derived
    class Manager that adds an attribute department. Write a program to input
    and display details of a manager using inheritance.
}

\codelabel

\begin{code}
#include <iostream>
#include <string>
using namespace std;

class Employee {
protected:
    string name;
    double salary;

public:
    void inputEmployee() {
        cout << "Enter name: ";
        getline(cin, name);
        cout << "Enter salary: ";
        cin >> salary;
        cin.ignore();
    }

    void displayEmployee() const {
        cout << "Name: " << name << endl;
        cout << "Salary: " << salary << endl;
    }
};

class Manager : public Employee {
    string department;

public:
    void input() {
        inputEmployee();
        cout << "Enter department: ";
        getline(cin, department);
    }

    void display() const {
        displayEmployee();
        cout << "Department: " << department << endl;
    }
};

int main() {
    Manager manager;
    manager.input();

    cout << "\nManager Details" << endl;
    manager.display();

    return 0;
}
\end{code}

\outputlabel
\outputimage[0.50\linewidth]{images/OOP/inheritance1.png}

\program{
    2. Design a class hierarchy with a base class Shape having a pure virtual
    function area(). Derive Rectangle and Circle from Shape and override area()
    in both. Demonstrate polymorphism by storing different shapes in an array
    of base class pointers and displaying their areas.
}

\codelabel

\begin{code}
#include <iostream>
using namespace std;

class Shape {
public:
    virtual double area() const = 0;
    virtual ~Shape() {}
};

class Rectangle : public Shape {
    double length, width;

public:
    Rectangle(double l, double w) {
        length = l;
        width = w;
    }

    double area() const override {
        return length * width;
    }
};

class Circle : public Shape {
    double radius;

public:
    Circle(double r) {
        radius = r;
    }

    double area() const override {
        return 3.14159 * radius * radius;
    }
};

int main() {
    Rectangle rectangle(10, 5);
    Circle circle(7);

    Shape* shapes[2] = {&rectangle, &circle};

    for (int i = 0; i < 2; i++)
        cout << "Area: " << shapes[i]->area() << endl;

    return 0;
}
\end{code}

\outputlabel
\outputimage[0.50\linewidth]{images/OOP/inheritance2.png}

\program{
    3. Build a class Vehicle with attributes make and model, and derive two
    classes Car and Bike. Car should have an attribute numDoors and Bike should
    have hasGear. Use constructors to initialize all data members and display
    details using overridden functions.
}

\codelabel

\begin{code}
#include <iostream>
#include <string>
using namespace std;

class Vehicle {
protected:
    string make, model;

public:
    Vehicle(string m, string mo) {
        make = m;
        model = mo;
    }

    virtual void display() const {
        cout << "Make: " << make << endl;
        cout << "Model: " << model << endl;
    }

    virtual ~Vehicle() {}
};

class Car : public Vehicle {
    int numDoors;

public:
    Car(string m, string mo, int doors)
        : Vehicle(m, mo), numDoors(doors) {}

    void display() const override {
        cout << "Car" << endl;
        Vehicle::display();
        cout << "Number of doors: " << numDoors << endl;
    }
};

class Bike : public Vehicle {
    bool hasGear;

public:
    Bike(string m, string mo, bool gear)
        : Vehicle(m, mo), hasGear(gear) {}

    void display() const override {
        cout << "Bike" << endl;
        Vehicle::display();
        cout << "Has gear: " << (hasGear ? "Yes" : "No") << endl;
    }
};

int main() {
    Car car("Toyota", "Corolla", 4);
    Bike bike("Honda", "Shine", true);

    car.display();
    cout << endl;
    bike.display();

    return 0;
}
\end{code}

\outputlabel
\outputimage[0.50\linewidth]{images/OOP/inheritance3.png}

\program{
    4. Implement a base class Person with name and age. Derive Student with
    rollNo and marks, and Teacher with subject and salary. Use base class
    pointers to call display() functions for objects of both derived classes
    and demonstrate runtime polymorphism.
}

\codelabel

\begin{code}
#include <iostream>
#include <string>
using namespace std;

class Person {
protected:
    string name;
    int age;

public:
    Person(string n, int a) : name(n), age(a) {}

    virtual void display() const {
        cout << "Name: " << name << endl;
        cout << "Age: " << age << endl;
    }

    virtual ~Person() {}
};

class Student : public Person {
    int rollNo;
    double marks;

public:
    Student(string n, int a, int r, double m)
        : Person(n, a), rollNo(r), marks(m) {}

    void display() const override {
        cout << "Student" << endl;
        Person::display();
        cout << "Roll Number: " << rollNo << endl;
        cout << "Marks: " << marks << endl;
    }
};

class Teacher : public Person {
    string subject;
    double salary;

public:
    Teacher(string n, int a, string s, double sal)
        : Person(n, a), subject(s), salary(sal) {}

    void display() const override {
        cout << "Teacher" << endl;
        Person::display();
        cout << "Subject: " << subject << endl;
        cout << "Salary: " << salary << endl;
    }
};

int main() {
    Student student("Rahul", 20, 101, 88.5);
    Teacher teacher("Mr. Sharma", 40, "C++", 50000);

    Person* people[2] = {&student, &teacher};

    for (int i = 0; i < 2; i++) {
        people[i]->display();
        cout << endl;
    }

    return 0;
}
\end{code}

\outputlabel
\outputimage[0.50\linewidth]{images/OOP/inheritance4.png}

\program{
    5. Create a class Account with data members accountNumber and balance.
    Derive two classes SavingsAccount and CurrentAccount which implement
    their own withdraw() functions with different rules. Demonstrate how
    function overriding works in this scenario.
}

\codelabel

\begin{code}
#include <iostream>
using namespace std;

class Account {
protected:
    int accountNumber;
    double balance;

public:
    Account(int number, double amount)
        : accountNumber(number), balance(amount) {}

    virtual void withdraw(double amount) = 0;

    void display() const {
        cout << "Account Number: " << accountNumber << endl;
        cout << "Balance: " << balance << endl;
    }

    virtual ~Account() {}
};

class SavingsAccount : public Account {
public:
    SavingsAccount(int number, double amount)
        : Account(number, amount) {}

    void withdraw(double amount) override {
        if (balance - amount >= 500) {
            balance -= amount;
            cout << "Savings withdrawal successful" << endl;
        } else {
            cout << "Savings account must maintain Rs. 500" << endl;
        }
    }
};

class CurrentAccount : public Account {
public:
    CurrentAccount(int number, double amount)
        : Account(number, amount) {}

    void withdraw(double amount) override {
        if (amount <= balance + 1000) {
            balance -= amount;
            cout << "Current withdrawal successful" << endl;
        } else {
            cout << "Overdraft limit exceeded" << endl;
        }
    }
};

int main() {
    SavingsAccount savings(101, 5000);
    CurrentAccount current(102, 5000);

    savings.withdraw(4700);
    current.withdraw(5500);

    cout << endl << "Savings Account" << endl;
    savings.display();

    cout << endl << "Current Account" << endl;
    current.display();

    return 0;
}
\end{code}

\outputlabel
\outputimage[0.50\linewidth]{images/OOP/inheritance5.png}

\program{
    6. Write a program that implements multilevel inheritance:
    Animal $\rightarrow$ Mammal $\rightarrow$ Dog. Each class should have a
    function describe() that prints information about the class. Show how
    the Dog class inherits features from both Animal and Mammal.
}

\codelabel

\begin{code}
#include <iostream>
using namespace std;

class Animal {
public:
    void describe() const {
        cout << "Animal: This is an animal." << endl;
    }
};

class Mammal : public Animal {
public:
    void describeMammal() const {
        cout << "Mammal: This is a mammal."
             << endl;
    }
};

class Dog : public Mammal {
public:
    void describeDog() const {
        cout << "Dog: This is a dog."
             << endl;
    }
};

int main() {
    Dog dog;

    dog.describe();
    dog.describeMammal();
    dog.describeDog();

    return 0;
}
\end{code}

\outputlabel
\outputimage[0.50\linewidth]{images/OOP/inheritance6.png}

\program{
    7. Implement multiple inheritance using two base classes Academic and
    Sports that both have marks as attributes, and a derived class Result that
    calculates the total score of a student by combining both. Handle ambiguity
    using scope resolution.
}

\codelabel

\begin{code}
#include <iostream>
using namespace std;

class Academic {
protected:
    int marks;

public:
    Academic(int m) : marks(m) {}

    void display() const {
        cout << "Academic Marks: " << marks << endl;
    }
};

class Sports {
protected:
    int marks;

public:
    Sports(int m) : marks(m) {}

    void display() const {
        cout << "Sports Marks: " << marks << endl;
    }
};

class Result : public Academic, public Sports {
public:
    Result(int academicMarks, int sportsMarks)
        : Academic(academicMarks), Sports(sportsMarks) {}

    void displayResult() const {
        Academic::display();
        Sports::display();
        cout << "Total Marks: " << Academic::marks + Sports::marks
             << endl;
    }
};

int main() {
    Result result(85, 15);
    result.displayResult();

    return 0;
}
\end{code}

\outputlabel
\outputimage[0.50\linewidth]{images/OOP/inheritance7.png}

\program{
    8. Write a program to illustrate protected access specifier. Make a base
    class Base with protected data, and derive a class Derived that modifies
    and displays this data through its own member functions.
}

\codelabel

\begin{code}
#include <iostream>
using namespace std;

class Base {
protected:
    int data;

public:
    Base() {
        data = 10;
    }
};

class Derived : public Base {
public:
    void modify() {
        data += 20;
    }

    void display() const {
        cout << "Data: " << data << endl;
    }
};

int main() {
    Derived object;

    object.modify();
    object.display();

    return 0;
}
\end{code}

\outputlabel
\outputimage[0.50\linewidth]{images/OOP/inheritance8.png}

\program{
    9. Implement a class hierarchy to represent Employee $\rightarrow$ Engineer
    $\rightarrow$ SoftwareEngineer. Use virtual inheritance to prevent ambiguity
    if SoftwareEngineer is also derived from another branch like
    Employee $\rightarrow$ Manager. Demonstrate how virtual inheritance resolves
    the diamond problem.
}

\codelabel

\begin{code}
#include <iostream>
#include <string>
using namespace std;

class Employee {
protected:
    string name;

public:
    Employee(string n) : name(n) {}

    void displayEmployee() const {
        cout << "Employee Name: " << name << endl;
    }
};

class Engineer : virtual public Employee {
public:
    Engineer(string n) : Employee(n) {}

    void displayEngineer() const {
        cout << "Engineer branch" << endl;
    }
};

class Manager : virtual public Employee {
public:
    Manager(string n) : Employee(n) {}

    void displayManager() const {
        cout << "Manager branch" << endl;
    }
};

class SoftwareEngineer : public Engineer, public Manager {
public:
    SoftwareEngineer(string n)
        : Employee(n), Engineer(n), Manager(n) {}

    void display() const {
        displayEmployee();
        displayEngineer();
        displayManager();
        cout << "SoftwareEngineer: Single Employee object shared."
             << endl;
    }
};

int main() {
    SoftwareEngineer softwareEngineer("Amit");
    softwareEngineer.display();

    return 0;
}
\end{code}

\outputlabel
\outputimage[0.50\linewidth]{images/OOP/inheritance9.png}
